\section{Related work and the design question}
\label{sec:related-work}

\subsection{Differential transfer functions and impedance balance}
Differential measurement requires both a useful differential response and a small common-to-differential response. Pall\`as-Areny and Webster derived differential-amplifier models with finite differential and common-mode gains and checked the predicted amplitude and phase experimentally \citep{pallasareny1991cmrr}. The companion analysis described cascaded differential stages using four transfer functions and stated the assumption that a stage does not load the preceding stage \citep{pallasareny1991cascaded}. For passive ladders connected directly to a receiver, the receiver impedance enters the branch transfer function. Loading therefore belongs in the matching equations from the outset.

The complex gain carries information needed to compare cancellation conditions. Koide used the complex differential-to-common-mode gain ratio to explain why deliberate imbalance can compensate an active device's common-mode error \citep{koide1997redef}. Gasulla and co-workers showed that common-mode gain and common-to-differential gain are needed to describe transient behavior and stability in fully differential circuits \citep{gasulla2002cmresp}. Hornero and co-workers subsequently analyzed how branch mismatch creates frequency-dependent common-mode responses, including band-pass behavior and phase shifts \citep{hornero2021cm}. The complex-response analyses motivate reporting the separate complex responses over frequency. The magnitude of a rejection ratio alone cannot identify which branch coefficients must agree.

\subsection{Topology, termination, and component matching}
Casas and Pall\`as-Areny analyzed uncoupled and coupled analog differential filters, including higher-order extensions \citep{casas1996basics}. Matching the transfers of uncoupled branches provides cancellation; changing the interconnection can instead produce cancellation through circuit topology. The floating common-point construction establishes an important reference case for topology-based rejection. The circuit considered here consists of two grounded two-stage ladders whose output terminals may carry equal resistive loads and an output-bridging capacitor. The component conditions must consequently be derived for those connections.

Impedance balance also determines the behavior of practical balanced interfaces. Whitlock and Floru showed how source and receiver common-mode impedance imbalances interact and described an input circuit with high common-mode impedance \citep{whitlock2004balancedinput}. The relevant design question for the present ladder is whether two branches with identical open-circuit voltage transfer continue to agree after termination. An open-circuit transfer fixes a voltage ratio, while the output impedance fixes how an attached load changes that ratio. Retaining both functions makes the termination dependence explicit and permits a direct comparison between open-circuit matching and matching at specified finite loads.

\subsection{Network realization and tolerance analysis}
Guillemin addressed RC-network synthesis; Fialkow and Gerst studied RC-ladder transfer functions and general two-terminal-pair RC networks \citep{guillemin1949rcsynthesis,fialkow1951rcladder,fialkow1952rcnetworks}. Wang and co-workers considered passive two-port RC-ladder realization with a specified gain \citep{wang2017rcladder}. The cited synthesis works establish the synthesis context for recovering positive components from transfer coefficients. The present calculation also retains the output-impedance numerator and asks how specified terminations change the set of components that match a fixed reference branch. \emph{Draft note: a proposition-level comparison with the four synthesis papers remains to be completed.}

Component realization and sensitivity must be considered together. Moschytz discussed the trade-offs among sensitivity, component spread, and component tolerance in active-filter design \citep{moschytz0radeffs}. Director and Rohrer developed a generalized adjoint-network framework for obtaining network sensitivities \citep{director1969the}. Endres and Kirkpatrick examined how differential loop-filter mismatch introduces additional transients in fast-settling phase-locked loops \citep{endres0ensitivi}. The sensitivity and transient studies support separating an exact nominal cancellation condition from the effects of component perturbations and from a settling-time specification. Each question requires a specified response function and evaluation criterion.

\subsection{Scope of the present contribution}
The present study derives the complete coefficient conditions for equality of the loaded transfers of two positive, grounded, two-stage RC branches. For a fixed reference branch, it constructs the admissible component family for an open circuit and for one positive resistive termination, and proves component uniqueness when equality holds at two distinct positive resistive terminations. It then examines how a capacitor connected only between the two outputs changes the separate common-to-differential and differential-to-differential gains. Analytical constructions are checked against native AC-simulation outputs.

The cancellation statements concern deterministic linear transfer functions. Serrano-Finetti and co-workers showed that common-mode electronic noise requires a source and cross-correlation model in addition to a deterministic common-mode rejection ratio \citep{serranofinetti2019cmnoise}. A noise-performance claim would therefore require the locations and correlations of the relevant noise sources. The present results supply the nominal network conditions and the termination dependence needed before that further calculation.
